\documentclass{article}
\usepackage{amsmath,amssymb,amsthm,mathtools,enumitem}
\usepackage[vocab=unitsquare]{leanify}
\newcommand{\E}{\mathbb{E}}
\newcommand{\CC}{\mathcal{C}}
\newcommand{\Gam}{\Gamma}
\newcommand{\R}{\mathbb{R}}
\newcommand{\de}{\,\mathrm{d}}
\newtheorem{theorem}{Theorem}
\newtheorem{lemma}[theorem]{Lemma}
\newtheorem{corollary}[theorem]{Corollary}
\begin{document}

A coupling is a probability measure on the unit square with uniform marginals: the law of a pair $(X,Z)$ of uniform random variables.
An expectation with a coupling as its subscript is the integral of a function of the two coordinates, which are written $X$ and $Z$.

Let $\CC$ denote the class of bivariate copulas.
Let $\Gam$ denote the set of couplings, that is, probability measures $\gamma$ on $[0,1]^2$ with uniform marginals.
For a copula $C$, Spearman's rho is
\begin{equation}
  \rho(C)=12\int_0^1\!\int_0^1 C(u,v)\de u\de v-3.
\end{equation}
For a copula $C$, we write $C^\top(u,v)\coloneqq C(v,u)$ for the transpose of $C$.
For a copula $C$, the survival copula of $C$ is $\widehat C(u,v)\coloneqq u+v-1+C(1-u,1-v)$.
For a Borel measurable $T:[0,1]\to[0,1]$, let $\gamma_T$ denote the law of $(X,T(X))$ for $X$ uniformly distributed on $[0,1]$.
For a Borel measurable $T:[0,1]\to[0,1]$, let $\gamma^T$ denote the law of $(T(Z),Z)$ for $Z$ uniformly distributed on $[0,1]$.

\begin{lemma}[Coupling representation]
\label{lem:coupling}
Let $C\in\CC$, and let $\gamma$ be the coupling associated with $C$.
Then
\begin{equation}
  \begin{aligned}
  \rho(C)&=12\,\E_\gamma[XZ]-3,
  &\rho(C^\top)&=12\,\E_\gamma[ZX]-3.
  \end{aligned}
\end{equation}
Consequently, $\rho$ is affine on $\CC$.
\end{lemma}

A family of maps of the interval gives a family of copulas, and a family may get one more value.

For $w\in[0,1]$, define the shuffle map $T_w:[0,1]\to[0,1]$ by
\begin{equation}
  T_w(x)\coloneqq
  \begin{cases}
    1-x, & 0\le x\le1-w,\\
    x+w-1, & 1-w<x\le1.
  \end{cases}
\end{equation}
For $w\in[0,1]$, let $A_w$ denote the copula associated with $\gamma_{T_w}$.
For $b\in(0,\frac12)$, let $B_b$ denote the copula associated with $\gamma^{T_b}$.
We therefore set $B_{1/2}\coloneqq A_{1/2}$.
For $b\in(0,\frac12)$ and $x\in(1-b,1]$, set
\begin{equation}
  z_-(x)\coloneqq\frac{x+b-1}{2(1-b)},
  \qquad
  z_+(x)\coloneqq\frac{1-b-(1-2b)x}{2(1-b)}.
\end{equation}

\begin{lemma}
\label{lem:values}
Let $w\in[0,1]$ and $b\in(0,\frac12]$.
Then $\rho(A_w)=2w^3-1$ and $\rho(B_b)\le1$.
Moreover, $w\mapsto\rho(A_w)$ increases bijectively from $[0,1]$ onto $[-1,1]$.
\end{lemma}

\begin{corollary}
\label{cor:sharp}
Every $C\in\CC$ satisfies $\rho(C)\le1$.
The maximum is attained only by $A_1$ and the minimum only by its survival copula, that is, for every $C\in\CC$, $\rho(C)=1$ if and only if $C=A_1$, and $\rho(C)=-1$ if and only if $C=\widehat{A_1}$.
Likewise, for every $w\in[\frac12,1]$, $\rho(A_w^\top)=2w^3-1$.
\end{corollary}

\begin{proof}
Let $r\in[-1,1]$, and let $C\in\CC$ with $\rho(C)=r$.
Let $\gamma\in\Gam$, and let $D$ be the copula associated with $\gamma$.
Then $\rho(D)\le1$ and $r\le1$.
\end{proof}

Potentials certify that a coupling is optimal.

\begin{lemma}[Dual certificates]
\label{lem:dual-certificate}
Let $c:[0,1]^2\to\R$ be a continuous function, and let $\varphi,\psi:[0,1]\to\R$ be continuous functions.
For $x,z\in[0,1]$, set $\Delta(x,z)\coloneqq\varphi(x)+\psi(z)-c(x,z)$.
Suppose that $\Delta(x,z)\ge0$ for all $x,z\in[0,1]$, and let $K\coloneqq\{(x,z)\in[0,1]^2:\Delta(x,z)=0\}$ be their contact set.
\begin{enumerate}[label=(\alph*),leftmargin=2em]
\item
Every $\gamma\in\Gam$ satisfies
\[
  \E_\gamma[c(X,Z)]\le\int_0^1\varphi(t)\de t+\int_0^1\psi(t)\de t,
\]
with equality if and only if $\gamma(K)=1$.
In particular, the maximizers are the couplings concentrated on $K$ when there is one, that is, if there is $\mu\in\Gam$ such that $\mu(K)=1$, then $\sup_{\gamma\in\Gam}\E_\gamma[c(X,Z)]=\int_0^1\varphi(t)\de t+\int_0^1\psi(t)\de t$.
\item
Let $T:[0,1]\to[0,1]$ be a Borel measurable map, and let $N_1$ and $N_2$ be finite subsets of $[0,1]$.
If $K\subseteq\{(x,z)\in[0,1]^2:z=T(x)\}\cup(N_1\times[0,1])\cup([0,1]\times N_2)$, then every $\gamma\in\Gam$ with $\gamma(K)=1$ satisfies $\gamma=\gamma_T$, and $T$ is Lebesgue measure preserving.
\end{enumerate}
\end{lemma}

Two random variables of a result are a random pair: the result is about every joint law of such a pair.

\begin{lemma}[Rearrangement inequality with equality case]
\label{lem:rearrangement}
Let $Y$ be an integrable random variable with continuous distribution function $F$, and let $Z$ be uniformly distributed on $[0,1]$ and defined on the same probability space.
Then $\E[YZ]\le\E[Y\,F(Y)]$, with equality if and only if $Z=F(Y)$ almost surely.
\end{lemma}

\begin{proof}
Fix $t\in(0,1)$ and set $q_t\coloneqq\sup\{s\in\R:F(s)\le t\}$.
Then $F(q_t)=t$.
\end{proof}

\end{document}
